\documentclass[11pt,a4paper]{article}

% --- Packages ---
\usepackage[utf8]{inputenc}
\usepackage[T1]{fontenc}
\usepackage[margin=1in]{geometry}
\usepackage{graphicx}
\usepackage{booktabs}
\usepackage{hyperref}
\usepackage{enumitem}
\usepackage{xcolor}
\usepackage{titlesec}
\usepackage{caption}

% --- Styling ---
\definecolor{darkblue}{rgb}{0.0, 0.0, 0.5}
\hypersetup{
    colorlinks=true,
    linkcolor=darkblue,
    urlcolor=darkblue,
    citecolor=darkblue
}

\titleformat{\section}{\large\bfseries\color{darkblue}}{\thesection}{1em}{}[\titlerule]
\titleformat{\subsection}{\normalsize\bfseries}{\thesubsection}{1em}{}

% --- Document Info ---
\title{
    \vspace{-2cm}
    \huge \textbf{Business Insights \& Explainability Report} \\
    \vspace{0.5cm}
    \large Phase 4: Customer Churn Predictor
}
\author{Data Science Team}
\date{\today}

\begin{document}

\maketitle

\section{Executive Summary}

We have successfully deployed a predictive AI model designed to identify customers at risk of leaving our service (churning). The model uses a \textit{Logistic Regression} approach, which provides a transparent and reliable way to rank customers by risk.

\subsection{How reliable is the model?}
To ensure the business can trust these predictions, we tested the model on a "hidden" set of data it had never seen before. The results are strong:
\begin{itemize}
    \item \textbf{Capture Rate (Recall) $\approx 82\%$}: The model successfully identifies over 80\% of all customers who actually churn. This means we have a high probability of catching "at-risk" customers before they leave.
    \item \textbf{Overall Accuracy $\approx 74\%$}: The model correctly predicts whether a customer will stay or leave roughly 3 out of 4 times.
    \item \textbf{Ranking Quality (ROC-AUC) $\approx 0.81$}: The model is very effective at distinguishing between a "loyal" customer and a "high-risk" customer, allowing us to prioritize our marketing spend efficiently.
\end{itemize}

\section{What Drives Customer Churn?}

We used a technique called \textbf{SHAP}, which acts like a "magnifying glass" to show exactly why the AI makes a specific prediction. Instead of a black box, we can now see the exact contribution of each customer behavior.

\subsection{The Top 5 Risk Factors}
The AI identified five primary behaviors that signal whether a customer is likely to stay or leave:

\begin{enumerate}
    \item \textbf{Purchase Frequency (Strongest Signal):} This is the most critical indicator. Customers who shop more often are significantly less likely to leave. A sudden drop in how often a customer visits is the earliest warning sign of attrition.
    \item \textbf{Monthly Revenue Generation:} Customers who generate a consistent, high amount of revenue per month are generally more stable and loyal.
    \item \textbf{Total Lifetime Spend (Monetary Value):} High-spending customers tend to have a deeper relationship with the brand and a higher "switching cost," making them more likely to stay.
    \item \textbf{Customer Tenure:} The longer a customer has been with us, the more stable they become. New customers are the most volatile and are most likely to churn early.
    \item \textbf{Variety of Products (NumProducts):} Customers who buy a wider range of different products are "stickier." Those who only buy one type of item are more vulnerable to competitor offers.
\end{enumerate}

% --- PLACEHOLDER: Insert plot showing SHAP Bar Plot here ---
\begin{figure}[htbp]
    \centering
    %\includegraphics[width=0.7\textwidth]{phase4_shap_bar.png}
    \fbox{Insert: phase4\_shap\_bar.png}
    \caption{Global Feature Importance: Which factors most influence the AI's decision?}
\end{figure}

\newpage

\section{Understanding the AI: Real-World Cases}

To demonstrate how the model works in practice, we analyzed three specific customer scenarios:

\subsection{Case A: The Accurate Warning (Correct Positive)}
The AI predicted this customer would churn, and they did.
\textbf{Why?} The model detected a "perfect storm" of warning signs: the customer stopped shopping entirely (High Recency) and had a very low purchase frequency. The AI correctly recognized that the lack of recent engagement meant the customer had already mentally "checked out."

\subsection{Case B: The Loyal Asset (Correct Negative)}
The AI predicted this customer would stay, and they remained active.
\textbf{Why?} The model saw strong evidence of stability. The customer shops almost weekly and has been with us for years. Even if their individual order values are small, their \textit{habitual} behavior signals extreme loyalty.

\subsection{Case C: The "False Alarm" (Misclassification)}
The AI flagged this customer as high-risk, but they actually stayed.
\textbf{Why?} This is a "Seasonal Buyer." They make huge purchases only a few times a year. The AI saw a long gap between orders and mistook it for abandonment. This highlights the need for us to identify "bulk buyers" to avoid wasting discounts on them.

% --- PLACEHOLDER: Insert SHAP Waterfall Plots ---
\begin{figure}[htbp]
    \centering
    \begin{minipage}{0.3\textwidth}
        \centering
        %\includegraphics[width=\textwidth]{phase4_correct_positive_waterfall.png}
        \fbox{Case A Plot}
        \caption{Correct Churn}
    \end{minipage}
    \hfill
    \begin{minipage}{0.3\textwidth}
        \centering
        %\includegraphics[width=\textwidth]{phase4_correct_negative_waterfall.png}
        \fbox{Case B Plot}
        \caption{Correct Stay}
    \end{minipage}
    \hfill
    \begin{minipage}{0.3\textwidth}
        \centering
        %\includegraphics[width=\textwidth]{phase4_misclassification_waterfall.png}
        \fbox{Case C Plot}
        \caption{False Alarm}
    \end{minipage}
\end{figure}

\section{Customer Risk Segmentation}

We have divided the customer base into three "Action Tiers" based on their probability of churning.

\begin{table}[htbp]
    \centering
    \caption{Business Profile by Risk Level}
    \vspace{0.3cm}
    \small
    \begin{tabular}{lccccp{3cm}}
        \toprule
        \textbf{Risk Tier} & \textbf{Avg Recency} & \textbf{Avg Freq} & \textbf{Avg Total Spend} & \textbf{Avg Order Value} & \textbf{Business Profile} \\
        \midrule
        \textbf{Low Risk} & 39 days & 181.2 & £4,365 & £27 & \textbf{Champions:} Highly active, high value. \\
        \addlinespace
        \textbf{Medium Risk} & 114 days & 41.6 & £1,053 & £305 & \textbf{Drifting:} High value, but infrequent. \\
        \addlinespace
        \textbf{High Risk} & 164 days & 12.3 & £257 & £34 & \textbf{At-Risk:} Dormant and low value. \\
        \bottomrule
    \end{tabular}
\end{table}

% --- PLACEHOLDER: Insert segment profiles plot ---
\begin{figure}[htbp]
    \centering
    %\includegraphics[width=0.8\textwidth]{phase4_segment_profiles.png}
    \fbox{Insert: phase4\_segment\_profiles.png}
    \caption{Comparison of behavioral traits across risk tiers.}
\end{figure}

\newpage

\section{Action Plan \& Roadmap}

Based on the AI's evidence, we recommend the following four initiatives:

\begin{description}[leftmargin=0em, labelsep=0.5em]
    \item[\textbf{1. Immediate High-Risk Outreach}] \hfill \\
    \textbf{Action:} Directly contact the "High Risk" segment within 24 hours via email/phone. \\
    \textbf{Why:} This group shows the strongest warning signals. Fast intervention is the only way to prevent total loss. \\
    \textbf{Priority:} \textbf{Critical} (Immediate).

    \item[\textbf{ la. Medium-Risk Nurture Journey}] \hfill \\
    \textbf{Action:} Implement an automated "We Miss You" email sequence with personalized offers. \\
    \textbf{Why:} These customers aren't gone yet, but they are drifting. A gentle nudge can bring them back to "Low Risk." \\
    \textbf{Priority:} \textbf{High} (Next 30 days).

    \item[\textbf{3. New Customer "Safety Net" Onboarding}] \hfill \\
    \textbf{Action:} Create a high-touch onboarding program for the first 90 days of a customer's life. \\
    \textbf{Why:} The AI shows that new customers (low tenure) are the most volatile. Establishing habits early reduces churn. \\
    \textbf{Priority:} \textbf{Medium} (Q2).

    \item[\textbf{4. Regional Calibration Check}] \hfill \\
    \textbf{Action:} Review model performance specifically for non-UK customers. \\
    \textbf{Why:} Because our data is mostly UK-based, the AI might be slightly less accurate for international clients. \\
    \textbf{Priority:} \textbf{Medium} (Quarterly Review).
\end{description}

\section{Ethical Considerations}

To ensure the AI is used fairly and responsibly, we have identified two key areas of focus:

\begin{itemize}
    \item \textbf{Geographic Fairness:} Since the majority of our customers are in the UK, there is a risk the model "ignores" the nuances of international markets. We will monitor "Error Rates" by country to ensure non-UK customers aren't unfairly ignored or misclassified.
    \item \textbf{New Customer Bias:} Customers with very short histories can look "risky" simply because the AI doesn't have enough data on them. We will implement a "Confidence Flag" for new accounts so we don't send aggressive "win-back" offers to people who just joined.
\end{itemize}

\end{document}
